\subsubsection{Overlap sweep and acceptance margin}
Figure~\ref{fig:X5} sweeps the side overlap from 50 to 70\,\% with the strip count and spacing recomputed at each value and the block centred in the polygon. At zero crab the accepted area equals the full \SI{100}{ha} only for side overlaps up to about \ncovMinOverlap\,\% and again at isolated higher values; around the claimed \SI{58}{\percent} the five-ray rule fails at the block edge for several consecutive overlaps (failing values include \ncovFailWindow\,\%). The number of strips steps from \ncovStripsSixFive{} at 65\,\% to \ncovStripsFiveEight{} at 58\,\%, saving \ncovTrackSave\,km of track. Crab angle changes the picture because the rotated footprint widens the cross-track extent: with crab $\ge3^\circ$ the failing window shrinks and with $\ge6^\circ$ it disappears, i.e.\ the acceptance of the 58\,\% plan depends on the wind---a fragility that a planner must treat as a constraint, not an average. The base-to-height ratio between adjacent strips grows linearly as overlap decreases, from \nBHbase{} at 65\,\% to \nBHopt{} at 58\,\%, which strengthens the intersection geometry between strips but, as stated above, says nothing about tie-point density.
\begin{figure}[H]\centering\includegraphics[width=\textwidth]{X5_overlap_sweep}
\caption{(a) Geometrically accepted area versus side overlap for five camera crab angles; (b) strips needed to cover \SI{1}{km}; (c) base-to-height ratio and minimum number of covering images (divided by 10).}\label{fig:X5}\end{figure}
